\makeatletter
%\@mainmatterfalse % Щоб не додавало в зміст
\chapter*{Апробація та практичне використання}
\makeatother

Результати, викладені в навчально-практичному посібнику «Протидія шахрайству: Платформи оголошень та цифровi псевдопослуги», пройшли апробацію в межах науково-практичних та наукових заходів всеукраїнського й міжнародного рівнів. Участь у цих заходах забезпечила обговорення результатів в науковому середовищі та підтвердила практичну значущість запропонованих підходів до протидії цифровим загрозам.

Перший етап апробації відбувся у 2023 році під час ІІІ Всеукраїнської науково-практичної Інтернет-конференції здобувачів освіти та молодих вчених «Маркетинг майбутнього»~\cite{Evdokymov2023}. У межах конференції було опубліковано матеріали на тему захисту інформації в умовах масового розповсюдження онлайн-реклами, що дозволило проаналізувати сучасні механізми маніпуляції користувачами та інформаційні ризики, пов’язані з цифровими маркетинговими інструментами. 

Наступним важливим етапом стала апробація результатів у 2024 році під час V Міжнародної науково-практичної онлайн-конференції «Обдаровані діти – скарб нації!», що відбулася в межах міжнародної освітньої виставки «Інноватика в сучасній освіті»~\cite{Yevdokymov2024}. У межах заходу результати дослідження були представлені у форматі майстер-класу, а також опубліковані у матеріалах конференції. Було продемонстровано можливості використання цифрових платформ моделювання як інструменту формування критичного мислення, цифрової грамотності та навичок безпечної роботи з технологіями, що є ключовими компетентностями в умовах зростання кіберзагроз.

У 2025 році результати дослідження отримали подальшу апробацію під час науково-практичної конференції «Штучний інтелект і безпека», де було опубліковано матеріали, присвячені моделюванню поведінки користувачів в умовах кібератак через фейкові додатки~\cite{Yevdokymov2025AI}. Представлені напрацювання дозволили розкрити роль штучного інтелекту в аналізі поведінкових сценаріїв та підтвердили актуальність досліджуваної проблематики в контексті сучасних кіберзагроз.

Окремим напрямом практичного впровадження результатів дослідження стало створення та розвиток відкритого освітнього цифрового ресурсу — вебсайту та мобільного застосунку «CyberSimulator»~\cite{SimulatorESA}. Зазначена освітня платформа розроблена автором на основі досліджень, матеріали якої викладені у книзі, веб-ресурс орієнтований на підвищення рівня цифрової грамотності й обізнаності користувачів у сфері кібербезпеки, що підтверджує прикладну спрямованість і соціальну значущість результатів дослідження.

Завершальним і ключовим етапом апробації стала ІV Міжнародна науково-практична конференція «Дизайн, візуальне мистецтво та творчість: сучасні тенденції та технології» (2025 р.), у межах якої відбувся публічний науковий виступ з доповіддю, присвяченою розпізнаванню графічних підробок у цифровому середовищі~\cite{Evdokymov2025Design}. Обговорення результатів на міжнародному рівні підтвердило їх наукову новизну та практичну значущість для протидії візуальному шахрайству, фейковому контенту та цифровим маніпуляціям.

Сукупність зазначених наукових заходів засвідчує системну апробацію результатів дослідження, їх фахове обговорення та практичну релевантність, що підкреслює наукову обґрунтованість і прикладну цінність матеріалів книги.



\clearpage % Завершуємо анотацію approbation
